% texte.tex — f168v, vv. 509–550 (conversion initiale des travaux antérieurs)
\vers[n=509]{}{}{τέξειν.' τὸ γὰρ ἦτρον τῆς χύτρας ἐλάκτισεν·}{accoucher ». Car il avait donné des pieds contre la panse de la marmite.}
\vers[n=510]{}{}{χὠ μὲν γεγηθὼς ἀπέτρεχ᾽, ἡ δ᾽ ἐξέσπασεν}{Tandis qu'il s'éloigne tout joyeux, elle débarrasse}
\vers[n=511]{}{}{ἐκ τοῦ στόματος τοῦ παιδίου, τὸ δ᾽ ἀνέκραγεν.}{la bouche de l'enfant qui se met à crier.}
\vers[n=512]{}{}{εἶθ᾽ ἡ μιαρὰ γραῦς, ἣ ᾽φερεν τὸ παιδίον,}{La vieille canaille qui l'a apporté}
\vers[n=513]{}{}{θεῖ μειδιῶσα πρὸς τὸν ἄνδρα καὶ λέγει,}{court en souriant à l'époux, et lui dit :}
\vers[n=514]{}{}{`λέων λέων σοι γέγονεν, αὐτέκμαγμα σόν,}{« un lion, un lion t'est né, ton portrait moulé ;}
\vers[n=515]{}{}{τά τ᾽ ἄλλ᾽ ἁπαξάπαντα καὶ τὸ πόσθιον}{outre tout le reste absolument, son membre}
\vers[n=516]{}{}{τῷ σῷ προσόμοιον, στρεβλὸν ὥσπερ κύτταρον.'}{est pareil au tien et tortu comme une pigne ».}
\vers[n=517]{}{}{ταῦτ᾽ οὐ ποιοῦμεν τὰ κακά; νὴ τὴν Ἄρτεμιν}{Ne sont-ce pas là de nos tours? Si fait, par Artémis,}
\vers[n=518]{}{}{ἡμεῖς γε. κᾆτ᾽ Εὐριπίδῃ θυμούμεθα,}{c'en sont. Et nous nous fâchons contre Euripide,}
\vers[n=519]{}{}{οὐδὲν παθοῦσαι μεῖζον ἢ δεδράκαμεν;}{« sans avoir enduré par-delà nos mérites ».;}
\vers[n=520, indent=2em, findent=1.2em, metre={lyr.}]{ΧΟ.}{Chœur}{τουτὶ μέντοι θαυμαστόν,}{Une chose me confond :}
\vers[n=521, indent=2em, findent=1.2em]{}{}{ὁπόθεν ηὑρέθη τὸ χρῆμα,}{où a-t-on trouvé cette créature ?}
\vers[n=522, indent=2em, findent=1.2em]{}{}{χἤτις ἐξέθρεψε χώρα}{Quel pays a formé}
\vers[n=523, indent=2em, findent=1.2em]{}{}{τήνδε τὴν θρασεῖαν οὕτω.}{une telle effrontée ?}
\vers[n=524, indent=2em, findent=1.2em]{}{}{τάδε γὰρ εἰπεῖν τὴν πανοῦργον}{La coquine ! tenir ce langage}
\vers[n=525, indent=2em, findent=1.2em]{}{}{κατὰ τὸ φανερὸν ὧδ᾽ ἀναιδῶς}{au grand jour, avec cette impudence ?}
\vers[n=526, indent=2em, findent=1.2em]{}{}{οὐκ ἂν ᾠόμην ἐν ἡμῖν}{je n'aurais pas cru que parmi nous}
\vers[suite]{ΧΟ.}{Chœur}{οὐδὲ τολμῆσαί ποτ᾽ ἄν.}{elle l'eût jamais osé.}
\vers[n=528, indent=2em, findent=1.2em]{}{}{ἀλλὰ πᾶν γένοιτ᾽ ἂν ἤδη·}{Mais maintenant tout arrive,}
\vers[n=529, indent=2em, findent=1.2em]{}{}{τὴν παροιμίαν δ᾽ ἐπαινῶ}{et j'approuve l'antique proverbe,}
\vers[n=529a, indent=2em, findent=1.2em]{}{}{τὴν παλαιάν· ὑπὸ λίθῳ γὰρ}{quasiment sous chaque pierre,}
\vers[n=530, indent=2em, findent=1.2em]{}{}{παντί που χρὴ}{il faut regarder}
\vers[suite]{ΧΟ.}{Chœur}{μὴ δάκῃ ῥητωρ ἀθρεῖν.}{de crainte d'être mordu... par un orateur..}
\vers[n=531, metre={ia⁴ cat.}]{ΚΟ.}{Coryphée}{ἀλλ᾽ οὐ γάρ ἐστι τῶν ἀναισχύτων φύσει γυναικῶν}{Mais les femmes au naturel impudent,}
\vers[n=532]{}{}{οὐδὲν κάκιον εἰς ἅπαντα πλὴν ἄρ᾽ εἰ γυναῖκες.}{il n'est rien de pire à tous égards, si ce n'est des femmes.}
\vers[n=533]{ΜΙ.}{Mika}{οὔτοι μὰ τὴν Ἄγλαυρον ὦ γυναῖκες εὖ φρονεῖτε,}{Non, par Aglaure, ô femmes, vous n'êtes pas dans votre bon sens, mais vous êtes maléficiées, ou vous avez pris quelque autre mal,}
\vers[n=534]{}{}{ἀλλ᾽ ἢ πεφάρμαχθ᾽ ἢ κακόν τι μέγα πεπόνθατ᾽ ἄλλο,}{pour que vous laissiez cette peste nous braver toutes à outrance.}
\vers[n=535]{}{}{ταύτην ἐῶσαι τὴν φθόρον τοιαῦτα περιυβρίζειν}{\didasc{Elle se tourne vers le bureau de l'assemblée.} Ainsi qu'on se décide...}
\vers[n=536]{}{}{ἡμᾶς ἁπάσας. εἰ μὲν οὖν τις ἔστιν· εἰ δὲ μή, ἡμεῖς}{sinon nous-mêmes et nos esclaves}
\vers[n=537]{}{}{αὐταί γε καὶ τὰ δουλάρια τέφραν ποθὲν λαβοῦσαι}{irons chercher quelque part de la cendre,}
\vers[n=538]{}{}{ταύτης ἀποψιλώσομεν τὸν χοῖρον, ἵνα διδαχθῇ}{et lui épilerons le bas-ventre, afin que, femme, elle apprenne}
\vers[n=539]{}{}{γυνὴ γυναῖκας οὖσα μὴ κακῶς λέγειν τὸ λοιπόν.}{à ne plus médire des femmes à l'avenir.}
\vers[n=540]{ΜΝ.}{Mnésiloque}{μὴ δῆτα τόν γε χοῖρον ὦ γυναῖκες. εἰ γὰρ οὔσης}{Le bas-ventre ? Non pas, ô femmes. Car la franchise}
\vers[n=541]{}{}{παρρησίας κἀξὸν λέγειν ὅσαι πάρεσμεν ἀσταί,}{étant de droit et nous toutes citoyennes ci-présentes ayant congé de parler,}
\vers[n=542]{}{}{εἶτ᾽ εἶπον ἅ ᾽γίγνωσκον ὑπὲρ Εὐριπίδου δίκαια,}{si j'ai déduit les raisons que je savais en faveur d'Euripide,}
\vers[n=543]{}{}{διὰ τοῦτο τιλλομένην με δεῖ δοῦναι δίκην ὑφ᾽ ὑμῶν;}{est-ce un motif pour m'épiler, afin de m'en faire porter la peine ?}
\vers[n=544]{ΜΙ.}{Mika}{οὐ γάρ σε δεῖ δοῦναι δίκην; ἥτις μόνη τέτληκας}{Quoi! tu n'en porterais pas la peine ? toi qui seule osas}
\vers[n=545]{}{}{ὑπὲρ ἀνδρὸς ἀντειπεῖν, ὃς ἡμᾶς πολλὰ κακὰ δέδρακεν}{répliquer en faveur d'un homme qui nous a tant malmenées,}
\vers[n=546]{}{}{ἐπίτηδες εὑρίσκων λόγους, ὅπου γυνὴ πονηρὰ}{trouvant à plaisir des sujets où une femme se montra mauvaise,}
\vers[n=547]{}{}{ἐγένετο, Μελανίππας ποιῶν Φαίδρας τε· Πηνελόπην δὲ}{produisant des Mélanippes, des Phèdres ; mais jamais il ne produisit une Pénélope,}
\vers[n=548]{}{}{οὐπώποτ᾽ ἐποίησ᾽, ὅτι γυνὴ σώφρων ἔδοξεν εἶναι.}{vu qu'elle passait pour sage.}
\vers[n=549]{ΜΝ.}{Mnésiloque}{ἐγὼ γὰρ οἶδα ταἴτιον. Μίαν γὰρ οὐκ ἂν εἴποις}{J'en sais bien la raison. Tu n'en citerais pas une seule,}
\vers[n=550]{}{}{τῶν νῦν γυναικῶν Πηνελόπην, Φαίδρας δ᾽ ἁπαξαπάσας.}{parmi les femmes d'aujourd'hui, de Pénélope. Toutes Phèdres, sans exception.}
